% Part: first-order-logic
% Chapter: syntax-and-semantics
% Section: formation-sequences

\documentclass[../../../include/open-logic-section]{subfiles}

\begin{document}

\olfileid{pl}{syn}{fseq}

\olsection{متتاليات التكوين}

من أحسن طرق البحث أن تُعرّف !!{formula}s استقرائيًا، ثم تُثبت خواص
!!{formula}s بالاستقراء. ومع حسن هذه الطريقة وقلة مؤونتها، قد ينفع أن نتابع بناء
!!{formula}s من مبادئه خطوة خطوة. ولأجل هذا نأخذ الآن بمفهوم
\emph{متتالية التكوين}.

\begin{defn}[متتاليات التكوين للصيغ]
\ollabel{defn:fseq-frm}
تكون المتتالية المنتهية $\tuple{!A_0,\dotsc,!A_n}$ المكوّنة من سلاسل
رموز من اللغة~$\Lang L_0$ \emph{متتالية تكوين} لـ~$!A$ إذا كان
$!A \ident !A_n$، وكان، لكل $i \leq n$، إما أن تكون $!A_i$ صيغة ذرية،
وإما أن يوجد $j,k < i$ بحيث تتحقق إحدى الحالات الآتية:
\begin{enumerate}
    \tagitem{prvNot}{$!A_i \ident \lnot !A_j$.}{}%
    \tagitem{prvAnd}{$!A_i \ident (!A_j \land !A_k)$.}{}%
    \tagitem{prvOr}{$!A_i \ident (!A_j \lor !A_k)$.}{}%
    \tagitem{prvIf}{$!A_i \ident (!A_j \lif !A_k)$.}{}%
    \tagitem{prvIff}{$!A_i \ident (!A_j \liff !A_k)$.}{}%
\end{enumerate}
\end{defn}

\begin{ex}
\[
\tuple{
    \Obj p_0,
    \Obj p_1,
    (\Obj p_1 \land \Obj p_0),
    \lnot (\Obj p_1 \land \Obj p_0)
}
\]
هي متتالية تكوين لـ
$\lnot (\Obj p_1 \land \Obj p_0)$، وكذلك
\[
\tuple{
    \Obj p_0,
    \Obj p_1,
    \Obj p_0,
    (\Obj p_1 \land \Obj p_0),
    (\Obj p_0 \lif \Obj p_1),
    \lnot (\Obj p_1 \land \Obj p_0)
}.
\]
%
وكما يتبين من المثال الثاني، قد تحتوي متتاليات التكوين على «حشو»:
أي صيغ زائدة أو لا تسهم في البناء.
\end{ex}

\begin{prop}
\ollabel{prop:formed}
لكل !!{formula}~$!A$ في~$\Frm[L_0]$ متتالية تكوين.
\end{prop}

\begin{proof}
افترض أن $!A$ ذرية. عندئذ تكون المتتالية~$\tuple{!A}$ متتالية
تكوين لـ~$!A$.
%
وافترض الآن أن لـ~$!B$ و~$!C$ متتاليتا تكوين
$\tuple{!B_0,\dotsc,!B_n}$ و~$\tuple{!C_0,\dotsc,!C_m}$،
على الترتيب.
%
\begin{enumerate}
  \tagitem{prvNot}{إذا كان $!A \ident \lnot !B$،
      فتكون $\tuple{!B_0,\dotsc,!B_n,\lnot !B_n}$
      متتالية تكوين لـ~$!A$.}{}
  \tagitem{prvAnd}{إذا كان $!A \ident (!B \land !C)$،
      فتكون $\tuple{!B_0,\dotsc,!B_n,!C_0,\dotsc,!C_m,(!B_n \land !C_m)}$
      متتالية تكوين لـ~$!A$.}{}
  \tagitem{prvOr}{إذا كان $!A \ident (!B \lor !C)$،
      فتكون $\tuple{!B_0,\dotsc,!B_n,!C_0,\dotsc,!C_m,(!B_n \lor !C_m)}$
      متتالية تكوين لـ~$!A$.}{}
  \tagitem{prvIf}{إذا كان $!A \ident (!B \lif !C)$،
      فتكون $\tuple{!B_0,\dotsc,!B_n,!C_0,\dotsc,!C_m,(!B_n \lif !C_m)}$
      متتالية تكوين لـ~$!A$.}{}
  \tagitem{prvIff}{إذا كان $!A \ident (!B \liff !C)$،
      فتكون $\tuple{!B_0,\dotsc,!B_n,!C_0,\dotsc,!C_m,(!B_n \liff !C_m)}$
      متتالية تكوين لـ~$!A$.}{}
\end{enumerate}
وبحسب مبدأ الاستقراء على !!{formula}s، لكل
!!{formula} متتالية تكوين.
\end{proof}

وأما العكس فيصح أيضًا، وإثباته يبيّن أن طريقتي تعريف الصيغ متكافئتان، لا تختلف نتائجهما. ومن ثمرته أن نثبت مبرهنات الصيغ بالاستقراء العادي على أطوال متتاليات تكوينها.

\begin{lem}
\ollabel{lem:fseq-init}
افترض أن $\tuple{!A_0,\dotsc,!A_n}$ متتالية تكوين لـ~$!A_n$، وأن
$k \leq n$. عندئذ تكون $\tuple{!A_0,\dotsc,!A_k}$ متتالية تكوين
لـ~$!A_k$.
\end{lem}

\begin{proof}
تمرين.
\end{proof}

\begin{thm}
\ollabel{thm:fseq-frm-equiv}
$\Frm[L_0]$ هي مجموعة جميع سلاسل الرموز في اللغة~$\Lang L_0$ التي
لها متتالية تكوين.
\end{thm}

\begin{proof}
اجعل $F$ مجموعة سلاسل رموز اللغة~$\Lang L_0$ التي لها متتالية تكوين.
قد ثبت في \olref[pl][syn][fseq]{prop:formed} أن
$\Frm[L_0] \subseteq F$؛ فليس علينا إلا إثبات الاحتواء الآخر.

لتكن لـ~$!A$ متتالية تكوين $\tuple{!A_0,\dotsc,!A_n}$.
والمطلوب $!A \in \Frm[L_0]$؛ فنستقرئ استقراء قويًا على~$n$.
نفرض أن كل سلسلة ذات متتالية تكوين رقم حدّها الأخير
$m < n$ تقع في~$\Frm[L_0]$. وتعريف متتالية التكوين يقسم الأمر: إما أن $!A_n$ ذرية، وإما أن يوجد $j,k < n$ وتتحقق إحدى الصور الآتية:
\begin{enumerate}
    \tagitem{prvNot}{$!A_n \ident \lnot !A_j$.}{}%
    \tagitem{prvAnd}{$!A_n \ident (!A_j \land !A_k)$.}{}%
    \tagitem{prvOr}{$!A_n \ident (!A_j \lor !A_k)$.}{}%
    \tagitem{prvIf}{$!A_n \ident (!A_j \lif !A_k)$.}{}%
    \tagitem{prvIff}{$!A_n \ident (!A_j \liff !A_k)$.}{}%
\end{enumerate}
فإن كانت $!A_n$ ذرية، فقد حصل
$!A_n \in \Frm[L_0]$. وإلا فلنأخذ حالة $!A \ident
(!A_j \land !A_k)$. تقضي \olref[pl][syn][fseq]{lem:fseq-init} بأن
$\tuple{!A_0,\dotsc,!A_j}$ و~$\tuple{!A_0,\dotsc,!A_k}$ متتاليتا تكوين لـ~$!A_j$ و~$!A_k$ على الترتيب. وهما مقطعان ابتدائيان حقيقيان من متتالية~$!A$؛ فرقم الحد الأخير في كل منهما أقل من~$n$. وبفرض الاستقراء يقع~$!A_j$ و~$!A_k$ في
$\Frm[L_0]$. ومن تعريف !!a{formula} يلزم أن
$(!A_j \land !A_k)$ كذلك. وسائر الحالات تُحمل على هذا المسلك.
\end{proof}

\end{document}
